\documentclass[11pt]{article}
\usepackage[T1]{fontenc}
\usepackage[margin=1in]{geometry}
\usepackage{amsmath,amssymb,microtype}
\usepackage[colorlinks=true,urlcolor=blue,citecolor=blue]{hyperref}
\title{The later canonical-form answer behind arXiv:1507.00613}
\author{Literature-status note for \texttt{fa\_banach\_001}}
\date{11 August 2026}

\begin{document}
\maketitle

\noindent\textbf{Status.}
\texttt{literature\_already\_answered}: canonical representation of
isometric monoid isomorphisms for convex, bounded-below, $1$-Lipschitz
functions on Banach spaces.

\section*{The exact source question}

For a Banach space $X$, let 
\[
 \mathcal{CL}^{1}(X)=\{f:X\to\mathbb R: f\text{ is convex and bounded
 below, and }f\text{ is }1\text{-Lipschitz}\},
\]
equipped with infimal convolution and the bounded transform of the uniform
metric.  Bachir's 2014 paper asks whether every isometric monoid
isomorphism between such function monoids has the canonical composition
form induced by an isometric isomorphism of the underlying Banach spaces
\cite{B14}.  The 2015 target restates that uncertainty immediately before
its Katetov-monoid Banach--Stone theorem: it says that the canonical form is
not known for ``the set of all convex $1$-Lipschitz bounded from below
functions defined on Banach space'' and points explicitly to Problem 2 of
the 2014 paper \cite{B15}.

\section*{The later affirmative answer}

In 2016 Bachir proved a complete canonical-form theorem for the larger
nonnegative $1$-Lipschitz monoids on invariant metric groups.  In the
notation of that paper, if
\[
 \Phi:(Lip^1_+(X),\oplus,\rho)\longrightarrow
       (Lip^1_+(Y),\oplus,\rho)
\]
is an isometric monoid isomorphism, then there is an isometric group
isomorphism $T:\overline X\to\overline Y$ such that
\[
 \Phi(f)=(\overline f\circ T^{-1})|_Y
 \qquad(f\in Lip^1_+(X)).
\]
This is Theorem 1 of arXiv:1610.02825 \cite{B16}.  For complete Banach
spaces, the completions in that statement are just $X$ and $Y$.

Most decisively, Remark 1(2) of the same paper states that following this
proof and replacing ``$1$-Lipschitz function'' by ``$1$-Lipschitz and
convex function'' gives a positive answer to Problem 2 of the 2014 paper.
This is an explicit question-to-answer cross-reference, not an inference
from related results.

The proof mechanism also explains why convexity causes no new ambiguity.
The monoid units identify the Kuratowski functions and hence the underlying
isometry $T$; isometry of the monoids recovers order and constant
translations; every $1$-Lipschitz function has the envelope
\[
 f=\inf_{t\in X}\{f(t)+\delta_t\};
\]
and the isomorphism preserves this infimum.  Applying the same argument
inside the convex submonoid yields $\Phi(f)=f\circ T^{-1}$.  Since an
isometric group isomorphism between real Banach spaces fixes $0$, the
Mazur--Ulam theorem makes $T$ linear, so this is the canonical Banach-space
form requested in Problem 2.

\section*{Conclusion and scope}

The sole live ``we do not know'' statement in arXiv:1507.00613 is therefore
not a fresh open problem.  It was answered affirmatively in
arXiv:1610.02825, with an explicit citation to the original problem.  The
second scanner hit in the target is inside a commented-out TeX line about
classifying continuous monoid morphisms to $\mathbb R$; it is not a
published open statement and is excluded from the mathematical scope.

\begin{thebibliography}{9}
\bibitem{B14}
M.~Bachir,
\emph{The inf-convolution as a law of monoid. An analogue to the
Banach--Stone theorem},
J. Math. Anal. Appl. 420 (2014), 145--166,
\url{https://doi.org/10.1016/j.jmaa.2014.05.074}.

\bibitem{B15}
M.~Bachir,
\emph{Any law of group metric invariant is an inf-convolution},
arXiv:1507.00613,
\url{https://arxiv.org/abs/1507.00613}.

\bibitem{B16}
M.~Bachir,
\emph{On representations of isometric isomorphisms between some monoid of
functions},
arXiv:1610.02825,
\url{https://arxiv.org/abs/1610.02825}.
\end{thebibliography}

\end{document}
